gives the exact common-side split: a certified proper descendant, or the
exceptional adjacent lift over the canonical lower LOSS token.  Hence neither
short row is a pure reset.

\item \textbf{Long high return.}  The returned valuation is $v\ge7$.  With
\[
 u=Q_{v-1}^g(J(t)),\qquad c=Q_{v-3}^g(3J(t)),
\]
the guard includes $O(c,\alpha(c))$ and the finite return provenance.
Lemma~\ref{lem:high-return-pay}(b) first makes the strict pair replacement
$(u,c)\mapsto(p,b)$.  The canonical alternative then returns directly to an
obligation at the already lower token $p$.  Only the factor alternative makes
the second strict replacement $(p,b)\mapsto(q,y)$ and enters the exact
marked-tail frame. Thus every long high return is strict, but only its factor
branch creates a marked tail.

\item \textbf{Marked tail.}  Its guard records a positive odd $C$, a phase $g$,
and an exact tail length $D$, with $b=Q_D^g(C)$ WIN and the exact marked
factor frame
\[
 \{Q_3^g(J(b)),Q_2^g(J(b))\}\text{ containing a \DRAW}.
\]
The retained WIN/LOSS pair is orientation-sensitive:
\[
\begin{array}{c|cc}
D&q\text{ WIN}&y\text{ LOSS}\\ \hline
1&B(b)&A(b)\\
2&A(b)&B(b)\\
D\ge3&A(b)&B(b).
\end{array}
\]
Lemma~\ref{lem:short-tail} proves that $D=1$ is itself strict: its first
factor source is an actual child of the retained LOSS token.  At $D=2$ the first source is an exact lift over the retained LOSS token
$y$. Lemmas~\ref{lem:loss-anchored-lift} and~\ref{lem:D2-pure} keep this
row attached: no unrelated finite endpoint is promoted, equal-token routing
only decreases its exact tail/factor counter, and every restart pays a strict
child-token edge.  For every $D\ge3$,
Lemma~\ref{lem:marked-tail-pay} replaces the retained LOSS token $y$ by its
actual WIN child
\[
 r=B(q)=B(y)\quad(D=3),\qquad
 r=B(q)=A(y)\quad(D\ge4),
\]
so $h(r)<h(y)$ before any new task is installed.

\item \textbf{Attached lift.}  There are two semantically distinct attached
lift origins, and they are not conflated.  The $D=2$ row is the exact lift over the already retained LOSS token $y$ of
Lemma~\ref{lem:short-tail}. It is normalized by
Lemma~\ref{lem:loss-anchored-lift}: while $y$ is fixed an exact tail/factor
counter decreases, and before any genuine restart $y$ is replaced by an
actual WIN child or certified descendant. It is not sent through the $D=3$
terminal-barrier proof.

The post-payment $D=3$ row has an exact raw lift $Q_m^g(J(r))$, $m\ge2$,
over the lower token $r$. Lemma~\ref{lem:attached-lift} gives its complete
three-row arithmetic.  The stable row $m\ge4$ performs only
$(m,a)\mapsto(m-3,9a)$, and the terminal rows are paid by
Lemma~\ref{lem:terminal-pay}.  Corollary~\ref{cor:short-lift-closed} closes
this module without any appeal to the $D=2$ row.

\item \textbf{Attached B lift.}  Every B-selecting exponent-one lift has the
form
\[
 S=Q_1^d(J(r)),\qquad e=1-d.
\]
Lemma~\ref{lem:attached-b} gives either terminal source or
$\rho(B(S))<r$.  Thus the selected source is a strict nested-source exit;
it is not a free reseed at the numerical cursor $S$.

\item \textbf{Post-payment $D=3$ terminal barriers.}  The terminal exponent is
exactly one, two, or three inside the shortest marked lift of (S8).
Lemma~\ref{lem:terminal-pay} treats all three.  The $D=2$ LOSS-anchored entry never uses this terminal lemma; it is already
covered by Lemmas~\ref{lem:loss-anchored-lift} and~\ref{lem:D2-pure}.
Exponent one
first reconstructs a lower LOSS token and then has only the symbolic suffix
rows $\lambda=1,2,3$ or $\lambda\ge4$; exponent two and positive-recycle
exponent three reconstruct a strictly lower LOSS token; the sole
zero-recycle exponent-three row replaces the existing LOSS token by its
actual WIN child.  Thus every terminal return pays before re-entering an
obligation.
\end{enumerate}

These rows cover every exponent and every returned tail length.  Therefore
P2 is closed: an upward control return is either preceded by a strict
retained source/token edge or is a finite same-retained-data route whose
remaining exponent is explicitly decreasing.
\end{proposition}

\begin{proof}
The coverage is by the exact parameters used to construct the controls, not
by a numerical cutoff.  Boundary and loss-sibling entries carry the true
finite local tail/factor counter $\tau$, and marked-tail controls carry their
true constant-tail exponent.  Factor forks split at $r=1$ versus $r\ge2$.
High returns split at the only two short valuations $5,6$ versus the uniform
range $v\ge7$.  Attached lifts split at $m=1,2,3,\ge4$.  Terminal rows split by
exponent $1,2,3$, and the exponent-one suffix is split by its actual
alternating-suffix length $\lambda=1,2,3,\ge4$.  These partitions are
disjoint except where two exact descriptions meet at a boundary, and in such
overlaps the certified conclusion is the same.

The parent/child incidences and rank payments are precisely the conclusions
of Lemmas~\ref{lem:factor-fork-pay}--\ref{lem:terminal-pay}.  A factor
alternative whose ordinary source is itself the retained anchor/token uses
the generic lower factor fork (S4).  A factor alternative at a temporary
short-lift cursor does \emph{not} use (S4): it remains inside (S8), keeps the
incoming $q/y$ provenance, and follows the exact finite lift/terminal
transport of Lemma~\ref{lem:attached-lift}.  Thus a temporary numerical
cursor is never promoted to a retained source merely to make a factor lemma
applicable.  The generic fork either pays a source/token component or reaches
(S5)--(S6).  A long return (S6) pays the pair-token edge before the marked
source is installed.  Its marked source then has an actual tail length $D$.
Rows $D\ge3$ pay the LOSS-token edge of (S7) before any raw/signed exit.
The $D=1$ row is itself strict. The $D=2$ row is not a free pure reset: it
remains LOSS-anchored by Lemma~\ref{lem:loss-anchored-lift}; equal-token
routing only decreases its exact tail/factor counter, and every restart first
pays a child-token edge.  Thus following the arrows gives the
closed cycle

\[
  \text{obligation}\to\text{factor fork}\to\text{high return}
  \to\text{marked tail}\to\text{attached lift}\to\text{obligation},
\]

but every traversal contains a strict retained edge before a new marked tail
is installed.  Inside the already post-payment $D=3$ lift, equal-retained-data
routing can additionally decrease the displayed natural lift exponent before
returning to the obligation.  The $D=2$ LOSS-anchored route is finite at fixed token/counter data and
cannot restart without the strict child-token event of
Lemma~\ref{lem:loss-anchored-lift}. Terminal exits in the $D=3$ module
are (S10) and pay a strict token first.  Therefore no finite raw cursor is ever promoted to
$M$ or $N$ merely because its outcome is known, and there is no semantic
``other'' row or unranked reverse edge.
\end{proof}

\section{Well-foundedness and the no-DRAW theorem}\label{sec:global}

\begin{lemma}[Well-founded-fibre principle]\label{lem:fibre-principle}
Let a transition relation carry a value in a well-order.  Suppose every
transition weakly decreases that value and the relation restricted to every
fixed value is well-founded.  Then the full relation is well-founded.
\end{lemma}

\begin{proof}
An infinite transition path would give a nonincreasing sequence in a
well-order.  It can contain only finitely many strict decreases, so from some
point onward the retained value would be constant.  The remaining infinite
path would lie inside one fixed fibre, contrary to the hypothesis.
\end{proof}

\begin{proposition}[Fixed-retained-data routing is well-founded]
\label{prop:fixed-fibre}
Fix the complete retained tuple
\[
 (s,\eta,M,\zeta,N,a).
\]
The routing relation that preserves every one of these values is
well-founded.
\end{proposition}

\begin{proof}
In a fixed fibre no seed transition, token replacement or source decrease is
available.  Proposition~\ref{prop:P2-closed} then leaves only finite routing
rows.

A boundary or loss-sibling tail has a finite exact exponent and decreases it
until an obligation/factor boundary.  A canonical A-side run cannot pass four
consecutive WIN side states (Corollary~\ref{cor:no-four-win}); before that it
reaches a B phase or one of the strict/attached exits.  A valuation-two B
reset cannot occur as a pure reverse edge: Proposition~\ref{prop:P1-closed}
requires a seed, source edge or token edge, all impossible in the fixed
fibre.  A larger B valuation enters a factor row.

A lower factor fork either leaves by a strict source/token edge or reaches a
high return.  Every high return is strict in a retained token by
Lemma~\ref{lem:high-return-pay}, so no high-return cycle survives in the
fixed fibre.  Consequently a marked tail cannot be \emph{re-entered} in a
fixed token fibre.  The $D=1$ row is itself strict, and every $D\ge3$ marked
row is strict by Lemma~\ref{lem:marked-tail-pay}.  The $D=2$ row may occur as an initial attached control after a high return,
but Lemma~\ref{lem:loss-anchored-lift} leaves no equal-fibre restart: while
its token is fixed the exact tail/factor counter decreases, and any return to
a fresh generic control is preceded by a strict replacement of that LOSS
token.

The post-payment $D=3$ shortest-lift module is the only remaining local module
with an unbounded displayed tail counter.  By Lemma~\ref{lem:attached-lift}
and Corollary~\ref{cor:short-lift-closed}, its sole same-provenance recurrence
is $m\mapsto m-3$ for $m\ge4$. Between two such updates the control moves
forward through finite raw-factor/terminal stages; at $m=1,2,3$ it meets a
terminal barrier and pays a token. Thus a fixed-fibre path is a finite
concatenation of decreasing natural counters and finite acyclic routing
segments. It cannot be infinite.
\end{proof}

\begin{proposition}[Well-founded marked normalizer]\label{prop:normalizer}
The complete marked routing relation is well-founded under the retained
projection $\Pi$ of \eqref{eq:v6-rank} together with
Proposition~\ref{prop:fixed-fibre}.
\end{proposition}

\begin{proof}
A strict outer-source transition lowers $s$ and may reset every later
component.  At fixed $s$, the unique transition $\eta:1\to0$ is strict and
may install $M$ and reset later data.  Once $\eta=0$, every change of $M$ is
a certified proper-descendant replacement and strictly lowers $\Phi(M)$ by
Lemma~\ref{lem:token-rank}; it may reset $\zeta,N,a$.

Only after $\eta=0$, at fixed preceding components, $\zeta:1\to0$ is a
one-time strict inner-token seed.  Thereafter every change of $N$ is a
certified descendant replacement
and strictly lowers $\Psi(N)$; it may reset $a$.  With both token layers
fixed, every change of the retained numerical anchor $a$ is a separately
proved strict source decrease.  Hence every transition weakly decreases
$\Pi$, and every transition that changes retained data decreases it
strictly.  The equal-$\Pi$ fibre is well-founded by
Proposition~\ref{prop:fixed-fibre}.  Lemma~\ref{lem:fibre-principle} now
proves the claim.
\end{proof}

\begin{proof}[Proof of Theorem~\ref{thm:main}]
Assume that a DRAW exists and choose the least coefficient source $s$ among
DRAW positions.  Start with the corresponding outer fibre and $\eta=1$.

If $s=0$, Lemma~\ref{lem:zero-source} performs a finite normalization to a
boundary or loss-sibling entry.  If $s>0$, follow a continuing DRAW through
long constant tails until exponent one or two.  Proposition~\ref{prop:boundary-reduction}
then gives a typed entry or the factor-free exponent-two base state; the
latter is handled by Proposition~\ref{prop:base-entry}.  The factorful
exponent-one case is handled by Proposition~\ref{prop:factorful-entry}; the
reverse-factor lemma is never used at exponent one.  At no point is a
decrease of a factorful canonical coefficient treated as a source decrease.

The first exact finite token may appear in a loss-sibling diamond, in the
factor-free base entry, or in the first free obligation.  Whichever typed
entry produces it, Lemma~\ref{lem:seed-policy} installs it in $M$ on the
unique transition $\eta:1\to0$; no earlier use of $\zeta$ is allowed.  If a
free obligation is reached before any finite endpoint, Proposition~\ref{prop:P1-closed}
either strictly decreases its retained source or produces exactly such a
seedable token.  Once $\eta=0$, the first additional finite witness of an
attached inner task may consume $\zeta$ once; after that, every free
$A/B_2$ reset is strict in an already declared source/token component or is
an attached route over existing provenance.  Every factor/high-return/marked-
tail/terminal route is exhaustive by Proposition~\ref{prop:P2-closed}.
Therefore all subsequent outcome-compatible normalization takes place in the
well-founded marked relation of Proposition~\ref{prop:normalizer}.

Every macro is justified by exact parent/child incidences and preserves an
actual DRAW continuation separately from its finite tokens.  A macro may use
a reverse parent as a proof transformation; it is not asserted to be a
forward play sequence.  Whenever it follows play, it uses an actual child.
A DRAW has no LOSS child and has at least one DRAW child, so whenever the
macro requires a continuing game branch a DRAW child exists.  Repeating the finite macros from the hypothetical DRAW
would therefore create an infinite path in the well-founded marked relation,
a contradiction.

Hence the conjugated game has no DRAW positions.  Proposition~\ref{prop:first-normal}
returns this conclusion to the original odd-integer game, with terminal
conjugated state $0$ corresponding to original state $1$.  Thus every
positive odd start has finite remoteness and optimal play reaches $1$.
\end{proof}


\section{Verification boundary and reproducibility}\label{sec:verification}

The theorem above is a human mathematical claim.  The accompanying Python
program checks the exact arithmetic identities used in Sections 2--9,
including the raw-selector and second-selector identities, both permanent
negative regressions, the exceptional no-$B_2$ assertion, and the declared
finite pure-control order.  It also checks the constant-tail, fixed-tail,
base-entry, zero-source and exceptional-side identities inherited unchanged
from version 5.0 over large finite ranges.

The finite program is supporting verification, not an infinite proof.  In
particular it does not assign outcomes to all integers and it does not replace
the proof-height arguments in Sections~\ref{sec:P2closed}--\ref{sec:global}.
The local high-return and terminal formulas used in the semantic closure are
cross-referenced, for audit provenance only, to Sections 91--135 of the frozen
repository proof ledger at revision
\nolinkurl{9d12789de3dbd841851f3faff1a237018d43879d}.  All statements needed for
the theorem are reproduced and proved in this manuscript.  The globally
audited but invalid entry statement formerly appearing as Section 137 of that
ledger is not a dependency of this manuscript.

No end-to-end Lean formalization is claimed.  The purpose of the machine files
is to make transcription and arithmetic regressions easy to attack
independently.

\appendix
\section{Permanent adversarial regressions}

The following examples are normative tests for any later revision.

\paragraph{Coefficient/source separation.}
For $s=1$, $J(1)=5$, $a=3J(1)=15$ and $e=1$, the common boundary child is
$R(44)=22$.  Its canonical coefficient is $11<15$, but $11=J(3)$, so its
coefficient source has grown from $1$ to $3$.  Therefore no factorful step may
use ``smaller coefficient'' as a synonym for ``smaller source''.

\paragraph{Pre-streak anchor separation.}
Lemma~\ref{lem:anchor-stress} gives
$20\to30\to45\to68$ and $B(68)=25>20$.  The proof of
Proposition~\ref{prop:P1-closed} deliberately compares only exact parent/child
proof tokens or the immediate typed inner source; it never compares $25$ with
$20$.

\section{Closed semantic inventory}\label{app:inventory}

\begin{center}
\small
\begin{tabular}{@{}p{.22\linewidth}p{.31\linewidth}p{.37\linewidth}@{}}
\toprule
Constructor & Exhaustive rows & Rank effect before restart \\
\midrule
loss-sibling entry & fixed-tail; exponent-one predecessor; side/hidden parent &
finite $\tau$ countdown while retaining the exact witness \\
free obligation & $A$ canonical; $B_2$; factor exit &
Prop.~\ref{prop:P1-closed}: source, seed, strict token, or attached entry \\
attached obligation & ordinary; exceptional exact lift &
retains existing token; factor/lift lemmas, never reseeds \\
factor fork & $r=1$; $r\ge2$ &
strict source/token or same-token high return \\
high return & $v=5$; $v=6$; $v\ge7$ &
strict descendant; for $v\ge7$ strict two-token multiset descent \\
marked tail & $D=1$; $D=2$; $D\ge3$ &
$D=1$ strict LOSS-child; $D=2$ LOSS-anchored countdown/child-token payment; $D\ge3$ strict $y\mapsto r$ \\
short lift & $D=2$ coupled entry; post-payment $D=3$ with $m\ge2$ &
$D=2$ stays LOSS-anchored and pays before restart; $D=3$ decreases $m$ by $3$ or pays a token \\
terminal & zero-recycle $m=2$; exponent $1$; positive $2/3$; zero $3$ &
strict descendant or exact lower-LOSS attached frame before restart \\
\bottomrule
\end{tabular}
\end{center}

\section{Frozen local-proof provenance map}

For reproducibility, the following map records where the same local incidence
calculations appeared in the research ledger.  It is an audit provenance map,
not an additional proof assumption: the theorem above uses the statements and
proofs reproduced in this manuscript.  At frozen revision
\nolinkurl{9d12789de3dbd841851f3faff1a237018d43879d}, the corresponding ranges
are:
\begin{itemize}[leftmargin=*]
\item Sections 91--95: obligation and lower factor-fork source/token exits;
\item Sections 97--103: short and long high-return provenance and the
      two-token descent;
\item Sections 104--120: transport of the lower token through raw/signed
      factor exits and exact lifts;
\item Sections 121--128: terminal reconstruction and the three finite gate
      lengths;
\item Sections 129--135: multiset token rank, fibre principle, terminal entry,
      the two short marked-tail rows, and the universal $D\ge3$ LOSS-token
      replacement.
\end{itemize}
Section 136 is useful as a historical assembly, but the present proof instead
uses Propositions~\ref{prop:P1-closed}, \ref{prop:P2-closed}, and
\ref{prop:normalizer}.  Sections 137--138 are explicitly excluded because the
former contained the audited coefficient/source shortcut.

\begin{thebibliography}{9}
\bibitem{Guy1986}
R. K. Guy, John Isbell's Game of Beanstalk and John Conway's Game of
Beans-Don't-Talk, \emph{Mathematics Magazine} 59(5) (1986), 259--269.

\bibitem{Guy1996}
R. K. Guy, Unsolved Problems in Combinatorial Games, Problem 42, in
\emph{Games of No Chance}, MSRI Publications 29, Cambridge University Press,
1996.

\bibitem{Althofer2024}
I. Alth\"ofer, M. Hartisch, T. Zipproth, Analysis of a Collatz Game and Other
Variants of the $3n+1$ Problem, \emph{Advances in Computer Games} (2024),
123--132. DOI: \url{https://doi.org/10.1007/978-3-031-54968-7_11}.

\bibitem{AlthoferPage}
I. Alth\"ofer, Collatz Problems: The Two-Player $3n+-1$ Game,
\url{https://althofer.de/collatz-prizes.html}.

\bibitem{Repo}
G. Pochuev, Optimal $3n\pm1$ Game: Proof and Verification Repository,
\url{https://github.com/Grisha-Pochuev/3n-plus-minus-1-game}.

\bibitem{Zenodo}
G. Pochuev, Optimal Play in the Two-Player $3n\pm1$ Game, Zenodo,
\url{https://doi.org/10.5281/zenodo.21844684}, 2026.
\end{thebibliography}

\end{document}
